\documentclass{article}

% Style NeurIPS 2025, lấy theo Call for Papers chính thức:
% https://neurips.cc/Conferences/2025/CallForPapers
% Nộp ẩn danh: \usepackage{neurips_2025}
% Bản báo cáo/preprint: \usepackage[preprint]{neurips_2025}
% Camera-ready: \usepackage[final]{neurips_2025}
\usepackage[preprint]{neurips_2025}

\usepackage{iftex}
\ifPDFTeX
  \usepackage[utf8]{inputenc}
  \usepackage[T1,T5]{fontenc}
  \usepackage[vietnamese,english]{babel}
\else
  \usepackage{fontspec}
  \setmainfont{Times New Roman}
  \usepackage[vietnamese,english]{babel}
\fi
\usepackage{hyperref}
\usepackage{url}
\usepackage{booktabs}
\usepackage{amsfonts}
\usepackage{amsmath}
\usepackage{graphicx}
\usepackage{microtype}
\usepackage{xcolor}

\selectlanguage{vietnamese}

\makeatletter
\renewcommand{\@noticestring}{Bản thảo báo cáo.}
\makeatother

\renewenvironment{abstract}%
{%
  \vskip 0.075in%
  \centerline{\large\bfseries Tóm tắt}%
  \vspace{0.5ex}%
  \begin{quote}%
}
{%
  \par%
  \end{quote}%
  \vskip 1ex%
}

\title{Dự báo lai chống rò rỉ dữ liệu cho Revenue và COGS theo ngày}

\author{%
  Nhóm Datathon VinU \\
  VinUniversity Datathon \\
  \texttt{team@example.com} \\
}

\begin{document}

\maketitle

\begin{abstract}
Báo cáo này trình bày hệ thống dự báo lai cho hai đại lượng tài chính theo
ngày là Revenue và Cost of Goods Sold (COGS) trên horizon 548 ngày. Phương pháp
kết hợp các feature biết trước theo thời gian, mốc lịch và khuyến mãi, seasonal
anchor, nhiều nhánh XGBoost, blending bằng out-of-fold predictions, trọng số
theo horizon bucket, hiệu chỉnh sai số trung vị, Kalman residual correction có
gate, và một nhánh reconciliation cho COGS dựa trên tỷ lệ COGS/Revenue. Pipeline
được thiết kế để tránh rò rỉ dữ liệu: tập test chỉ cung cấp ngày và các biến
known-future, còn target lag được sinh từ lịch sử quá khứ hoặc từ prediction
recursive trước đó. Rolling-origin validation được dùng để mô phỏng đúng bài
toán dự báo mà không train trên nhãn tương lai.
\end{abstract}

\section{Giới thiệu}

Bài toán của cuộc thi là dự báo chuỗi thời gian tài chính dài hạn ở tần suất
ngày. Mô hình cần dự báo đồng thời Revenue và COGS cho giai đoạn submission dựa
trên lịch sử doanh số hằng ngày và metadata khuyến mãi nếu có. Bối cảnh này
khác các bài toán demand forecasting ngắn hạn, vì phần lớn horizon không thể dựa
vào các lag ngắn đã quan sát được. Do đó, lời giải cần cân bằng giữa ngoại suy
mùa vụ, feature known-future ổn định, mô hình gradient boosting trực tiếp và các
bước hậu xử lý bảo thủ.

Báo cáo này mô tả notebook cuối cùng của hướng V5 Hybrid. Các đóng góp chính
gồm: (i) feature store deterministic chỉ dùng calendar, event và promotion-plan
biết trước; (ii) rolling-origin validation theo đúng chiều thời gian; (iii) thư
viện candidate XGBoost đa dạng và được blend bằng out-of-fold predictions; (iv)
phân rã target COGS qua nhánh tỷ lệ COGS/Revenue; và (v) các diagnostics giúp
giải thích đóng góp của từng candidate và từng bước hậu xử lý.

\section{Phương pháp tiếp cận}

\subsection{Ý tưởng tổng quát}

Gọi $y^{R}_{t}$ là Revenue và $y^{C}_{t}$ là COGS tại ngày $t$. Hệ thống dự báo
Revenue trực tiếp, dự báo COGS trực tiếp, đồng thời dự báo thêm tỷ lệ chi phí
\begin{equation}
  r_t = \mathrm{clip}\left(\frac{y^{C}_{t}}{\max(y^{R}_{t}, 1)}, 0.01, 3.0\right).
\end{equation}
COGS cuối cùng được reconcile theo công thức
\begin{equation}
  \hat{y}^{C}_t =
  \alpha \hat{y}^{C,\mathrm{direct}}_t +
  (1-\alpha)\hat{y}^{R}_t \hat{r}_t,
\end{equation}
trong đó $\alpha \in [0,1]$ được chọn bằng MAE trên out-of-fold validation. Cách
làm này đưa ràng buộc kinh doanh giữa doanh thu và chi phí vào pipeline, nhưng
vẫn giữ một mô hình COGS trực tiếp để bắt các biến động chi phí không được giải
thích hoàn toàn bởi Revenue.

Ma trận feature được chia thành hai nhóm. Nhóm known-future gồm các biến từ
ngày tháng, Fourier terms, tháng/quý, khoảng cách tới Tết, double-day events,
Black Friday/Cyber Monday, payday windows và các thống kê từ kế hoạch khuyến
mãi. Nhóm past-only gồm lag target và rolling statistics. Trong validation và
test, feature lag được sinh theo kiểu recursive: sau khi mô hình dự báo ngày
$t$, prediction đó có thể trở thành lịch sử cho ngày $t+1$, nhưng actual của
validation hoặc test không bao giờ được đưa vào lịch sử recursive.

\subsection{Pipeline mô hình}

Bảng~\ref{tab:pipeline} tóm tắt pipeline thực thi trong notebook. Quy trình này
được chạy cho Revenue, COGS trực tiếp và COGS ratio; riêng COGS ratio dùng một
model library nhỏ hơn.

\begin{table}[h]
  \centering
  \caption{Pipeline V5 Hybrid được triển khai trong notebook.}
  \label{tab:pipeline}
  \small
  \begin{tabular}{p{0.16\linewidth}p{0.27\linewidth}p{0.45\linewidth}}
    \toprule
    Bước & Hàm/khối code & Vai trò \\
    \midrule
    Đọc dữ liệu & \texttt{load\_data} &
      Đọc sales, sample submission và promotions nếu có. \\
    Feature store & \texttt{build\_feature\_frame} &
      Tạo calendar, event và promotion features biết trước. \\
    Lag features & \texttt{make\_lag\_features} &
      Tạo lag và rolling statistics đã shift từ target quá khứ. \\
    Validation & \texttt{make\_splits} &
      Tạo rolling-origin expanding-window folds. \\
    Train target & \texttt{train\_target\_hybrid} &
      Fit anchors, XGBoost heads, blend, calibration và Kalman gate. \\
    Reconciliation & \texttt{run\_pipeline} &
      Kết hợp COGS trực tiếp và Revenue nhân với COGS ratio. \\
    Output & \texttt{run\_pipeline} &
      Ghi submission, OOF predictions, component forecasts và summary JSON. \\
    \bottomrule
  \end{tabular}
\end{table}

Model library cho Revenue và COGS trực tiếp gồm seasonal anchor và năm nhánh
XGBoost: hai mô hình all-history trên log-scale với squared-error, một mô hình
recent-history trên log-scale, một mô hình recent-history với absolute-error, và
một mô hình Tweedie. Một mô hình lag recursive có thể được bật để bổ sung tín
hiệu từ quá khứ gần. Nhánh COGS ratio dùng hai XGBoost head nhỏ hơn: một mô
hình all-history và một mô hình recent-history absolute-error. XGBoost được dùng
vì tree ensembles xử lý tốt các tương tác phi tuyến giữa calendar, event và
promotion features, đồng thời vẫn nhanh và dễ tái lập \citep{xgboost}.

Với mỗi target, các prediction OOF được blend bằng trọng số không âm, tối ưu
theo MAE thông qua các mẫu Dirichlet. Trọng số blend có thể được chuyên biệt
theo horizon bucket: ngày 1--30, 31--90, 91--180, 181--365 và 366+. Sau đó,
pipeline chọn multiplier và median residual bias để calibration. Cuối cùng,
damped Kalman residual projector chỉ được giữ nếu grid search trên OOF folds cải
thiện MAE; nếu không, strength của Kalman được đặt bằng 0.

\subsection{Thiết kế validation theo thời gian}

Cấu hình mặc định dùng validation horizon 548 ngày, step size 365 ngày và tối đa
3 folds. Mỗi fold là một expanding-window backtest: train indices luôn đứng
trước validation indices. Vì validation horizon dài hơn step size, các validation
windows liền kề có thể overlap; tuy nhiên, không fold nào train trên nhãn của
chính validation period hoặc trên bất kỳ nhãn tương lai nào của test period. Đây
là ràng buộc quan trọng nhất để cross-validation đúng cho bài toán chuỗi thời
gian.

\section{Thiết lập thực nghiệm}

Pipeline đọc \texttt{sales.csv}, \texttt{sample\_submission.csv} và dùng
\texttt{promotions.csv} nếu file này tồn tại. Ngày được chuẩn hóa, các dòng sales
trùng ngày bị loại bỏ, Revenue/COGS được ép kiểu số và clip không âm. File sample
được rút gọn chỉ còn cột \texttt{Date} trước khi tạo feature, nên target trong
template test không thể đi vào mô hình.

Metric nội bộ chính là MAE; RMSE và $R^2$ được in thêm để diagnostic. Notebook
ghi ba file OOF: \texttt{v5\_oof\_revenue.csv},
\texttt{v5\_oof\_cogs\_direct.csv} và \texttt{v5\_oof\_cogs\_ratio.csv}. Các
artifact này chứa actual values, prediction của từng candidate, global blend,
horizon-bucket blend, calibrated predictions và Kalman-gated output. Đây là cơ
sở thực nghiệm chính để so sánh từng head và từng bước hậu xử lý.

\section{Kết quả thực nghiệm}

\subsection{So sánh mô hình trên OOF}

Khi chạy notebook, mỗi target sẽ in một bảng metric. Bảng này so sánh seasonal
anchor, từng XGBoost candidate, global blend, horizon-bucket blend, calibrated
blend và Kalman-gated forecast. File summary JSON lưu lại global weights, bucket
weights, calibration parameters, Kalman strength và COGS reconciliation weight.
Nhờ vậy ensemble có thể được giải thích ở cấp mô hình: weight lớn cho thấy OOF
search tin candidate đó nhiều hơn cho target hoặc horizon bucket tương ứng; còn
Kalman strength bằng 0 nghĩa là residual correction đã bị validation gate loại.

Bảng~\ref{tab:diagnostics} liệt kê các diagnostics được tạo bởi một lần chạy đầy
đủ và cách đọc chúng trong báo cáo.

\begin{table}[h]
  \centering
  \caption{Các artifact thực nghiệm được tạo bởi notebook V5 Hybrid.}
  \label{tab:diagnostics}
  \small
  \begin{tabular}{p{0.39\linewidth}p{0.52\linewidth}}
    \toprule
    Artifact & Ý nghĩa \\
    \midrule
    \path{v5_summary.json} &
      Metrics, blend weights, calibration, Kalman gate và COGS alpha. \\
    \path{v5_oof_*.csv} &
      Actual và prediction theo fold để so sánh từng candidate. \\
    \path{v5_final_*_components.csv} &
      Forecast component trên test period trước reconciliation cuối. \\
    \path{submission_v5_hybrid_model_only.csv} &
      Forecast thuần từ mô hình, chưa anchor blend. \\
    \path{submission_v5_recommended.csv} &
      Submission được khuyến nghị, có thể đã anchor-blend sau forecast. \\
    \bottomrule
  \end{tabular}
\end{table}

\subsection{Kiểm tra phân phối submission}

Snapshot hiện tại của repository có file \texttt{submission\_v5\_recommended.csv}
gồm 548 dòng dự báo từ 2023-01-01 đến 2024-07-01. Bảng~\ref{tab:submission_summary}
tóm tắt file forecast này. Đây không phải validation score; đây là kiểm tra phân
phối output để xác nhận horizon, scale và non-negativity.

\begin{table}[h]
  \centering
  \caption{Thống kê tóm tắt của recommended submission đã sinh.}
  \label{tab:submission_summary}
  \begin{tabular}{lrrrrrr}
    \toprule
    Target & Mean & Std. & Min & Median & Max & Dòng âm \\
    \midrule
    Revenue & 3.309M & 1.660M & 0.806M & 3.103M & 10.339M & 0 \\
    COGS    & 3.064M & 1.540M & 0.649M & 2.838M &  9.496M & 0 \\
    \bottomrule
  \end{tabular}
\end{table}

\subsection{Giải thích mô hình}

Notebook cung cấp explainability ở cấp mô hình thông qua OOF component metrics,
blend weights, horizon-bucket weights, calibration parameters và quyết định
Kalman gate. Đây là lớp giải thích an toàn nhất vì được tính từ chính validation
predictions dùng để chọn ensemble cuối. Nếu cần giải thích feature-level, có thể
chạy SHAP trên XGBoost head được chọn như một bước post-hoc mà không thay đổi
forecast:
\begin{equation}
  I_j = \frac{1}{n}\sum_{i=1}^{n} |\phi_{ij}|,
\end{equation}
trong đó $\phi_{ij}$ là SHAP contribution của feature $j$ tại mẫu $i$. Các giá
trị $I_j$ sau khi xếp hạng có thể được báo cáo cùng các diagnostics ở cấp mô
hình.

\section{Hạn chế}

Phương pháp này được thiết kế bảo thủ và hướng tới leaderboard. Thứ nhất, mô
hình chỉ dùng chuỗi tổng hợp, chưa khai thác hierarchy theo sản phẩm, kênh hoặc
cửa hàng. Thứ hai, validation windows overlap giúp có thêm long-horizon
backtests nhưng không tạo ra các mẫu lỗi hoàn toàn độc lập. Thứ ba, recursive
lag forecasting có thể tích lũy sai số ở horizon dài, dù ensemble đã giảm rủi ro
bằng deterministic known-future heads và seasonal anchors. Cuối cùng, anchor
blend là chiến lược hậu xử lý submission, không phải cải thiện mô hình được
validate, nên được tách khỏi training pipeline.

\section{Kết luận}

V5 Hybrid kết hợp leakage-safe feature engineering, validation đúng chiều thời
gian, model library XGBoost đa dạng, OOF blending, calibration, Kalman residual
gating và COGS reconciliation có cấu trúc. Pipeline này có khả năng tái lập,
giải thích được ở cấp component và phù hợp với ràng buộc của bài toán dự báo tài
chính dài hạn theo ngày.

\section*{Lời cảm ơn}

Báo cáo này được chuẩn bị cho phần thuyết minh dự án tại VinUniversity Datathon.

\section*{References}

\bibliographystyle{plainnat}
\bibliography{references}


\appendix

\section{Chi tiết thực nghiệm bổ sung}

Notebook được thiết kế để chạy trực tiếp trên Kaggle hoặc môi trường local có
đủ các file dữ liệu của cuộc thi. Các tham số quan trọng gồm
\texttt{val\_days = 548}, \texttt{step\_days = 365},
\texttt{max\_folds = 3}, \texttt{n\_blend\_samples = 12000}, và
\texttt{seed = 42}. Khi cần smoke test nhanh, có thể bật
\texttt{RUN\_FAST\_LOCAL=True} để giảm số cây của
các mô hình XGBoost.

\section{Kết quả bổ sung}

Các bảng OOF chi tiết nằm trong các file \path{v5_oof_*.csv}. Các file
\path{v5_final_*_components.csv} cho phép kiểm tra từng component forecast
trong giai đoạn submission, bao gồm seasonal anchor, từng candidate XGBoost,
\texttt{blend\_bucket}, \texttt{calibrated} và \texttt{kalman\_gated}.

\section{Hạn chế và tác động rộng hơn}

Mô hình phục vụ mục tiêu dự báo vận hành và lập kế hoạch tài chính. Rủi ro chính
là forecast sai có thể dẫn đến lập kế hoạch tồn kho hoặc ngân sách chưa tối ưu.
Do đó, prediction nên được dùng cùng kiểm tra nghiệp vụ và không nên xem là một
quyết định tự động không cần giám sát.


% Call for Papers của NeurIPS 2025 yêu cầu checklist khi nộp thật.
% Khi cần nộp theo NeurIPS, bỏ comment các dòng dưới và điền checklist.tex.
% \newpage
% \input{checklist.tex}

\end{document}
